\begin{table}[!htbp]\centering
\caption{Main premium-enrollment outcomes: cubic-score OLS}
\label{tab:premium-cubic-income-main}
\small
\begin{tabular}{lrrrr}\hline\hline
 & Institution & Field & Either & Both \\ \hline
Math: average derivative (per SD) & 7.27 (0.18) & 12.75 (0.13) & 13.56 (0.11) & 6.47 (0.15) \\
Verbal: average derivative (per SD) & 4.62 (0.14) & 1.37 (0.13) & 4.05 (0.13) & 1.94 (0.09) \\
Income decile 1 & Reference & Reference & Reference & Reference \\
Income decile 2 vs. 1 & 0.73 (0.16) & 0.14 (0.31) & 0.49 (0.31) & 0.38 (0.13) \\
Income decile 3 vs. 1 & 1.79 (0.18) & 1.19 (0.32) & 2.06 (0.32) & 0.92 (0.14) \\
Income decile 4 vs. 1 & 1.90 (0.18) & 1.64 (0.32) & 2.46 (0.32) & 1.08 (0.14) \\
Income decile 5 vs. 1 & 3.27 (0.22) & 1.87 (0.33) & 3.51 (0.33) & 1.62 (0.16) \\
Income decile 6 vs. 1 & 3.78 (0.24) & 2.07 (0.34) & 3.93 (0.34) & 1.92 (0.17) \\
Income decile 7 vs. 1 & 4.89 (0.29) & 2.00 (0.35) & 4.42 (0.36) & 2.47 (0.19) \\
Income decile 8 vs. 1 & 6.11 (0.34) & 1.95 (0.36) & 5.13 (0.38) & 2.93 (0.22) \\
Income decile 9 vs. 1 & 9.67 (0.45) & 1.27 (0.37) & 6.06 (0.40) & 4.89 (0.29) \\
Income decile 10 vs. 1 & 25.88 (0.94) & -7.30 (0.59) & 9.47 (0.47) & 9.11 (0.48) \\
\hline
Outcome mean (\%) & 18.57 & 38.93 & 46.98 & 10.52 \\
Full R-squared & 0.2493 & 0.0978 & 0.1729 & 0.1783 \\
R-squared without income & 0.2180 & 0.0947 & 0.1706 & 0.1721 \\
Income share of explained variation (\%) & 12.5 & 3.2 & 1.3 & 3.5 \\
Income indicators: joint clustered p-value & $<$0.001 & $<$0.001 & $<$0.001 & $<$0.001 \\
Observations & 565,040 & 565,040 & 565,040 & 565,040 \\
School clusters & 3,372 & 3,372 & 3,372 & 3,372 \\
Admission-year indicators & Yes & Yes & Yes & Yes \\
\hline\hline\end{tabular}
\begin{minipage}{0.98\textwidth}\footnotesize
\textit{Notes:} Income-decile coefficients and score average derivatives are in percentage points; school-clustered CR1 SEs are in parentheses. Score controls include $M,V,M^2,V^2,MV,M^3,V^3,M^2V,MV^2$. Income decile 1 is the reference. Admission-year effects are included in both full and reduced models. Income share is $(R^2_{full}-R^2_{without\ income})/R^2_{full}$, expressed as a percentage; it measures incremental explanatory power after scores and year, not a causal contribution. The original OLS sample, including 5,105 admission-2026 records and all observed non-enrollees, is retained. Engineering includes manufacturing/construction; Medicine+ retains the six-program health-field definition.
\end{minipage}
\end{table}
